Über den Schnorchel ist die Lunge mit der Aussenluft verbunden, in ihr herrscht also der Luftdruck $p_\text{L}$. Auf den Brustkorb drückt aber das Wasser mit $p_\text{L} + \rho_\text{W}\,g\,h$. Die Differenz ist der Schweredruck des Wassers ($\rho_\text{W} = \rhoWO$, FoTaBe S.~187):
\begin{align}
 \Delta p &= \pueF \\
   &= \rhoW\cdot\gE\cdot\hT \\
   &= \pue \approx \result{\puebP{-3}{2}}
\end{align}
Schon dieser kleine Überdruck macht das Einatmen sehr anstrengend; deshalb sind Schnorchel kurz.
